% ==============================================================================
% SEZIONE 6: EVOLUZIONE DEL MIX PRODUTTIVO E CAPACITÀ INSTALLATA
% ==============================================================================

\section{Mix Produttivo}

\begin{frame}{Evoluzione del Mix Produttivo Globale}
    La strategia di Enel punta a un mix produttivo bilanciato, dove la crescita delle rinnovabili è supportata da tecnologie di stabilità e sistemi di accumulo.
    \vspace{0.1cm}
    
    \note{La capacità installata globale ha raggiunto quasi ottantasettemila megawatt, fondandosi su tre pilastri strategici.}

    \begin{block}{Ripartizione della Capacità Totale (86.986 MW)}
        \begin{itemize}
            \item \textbf{Renewables Core:} Idroelettrico (28.320 MW), Eolico (16.184 MW), Solare (13.059 MW).
            \pause
            \note{Il nucleo rinnovabile è guidato dall'idroelettrico, seguito da eolico e solare.}
            
            \item \textbf{Storage \& Support:} \textbf{BESS} (Accumulo a Batterie) per \textbf{3.441 MW}, affiancati da Nucleare (3.328 MW) e CCGT (12.420 MW).
            \pause
            \note{La stabilità è garantita dagli accumuli a batterie BESS, supportati da nucleare e impianti a ciclo combinato.}
            
            \item \textbf{Phase-out Fossili:} Il Carbone scende a \textbf{4.627 MW}, segnando il declino irreversibile della generazione termica tradizionale.
        \end{itemize}
    \end{block}
    \note{Il progressivo phase out dei combustibili fossili riduce il carbone a soli quattromila-seicentoventisette megawatt.}
\end{frame}


\begin{frame}{Specializzazione Geografica nelle Aree Tier 1}
    L'analisi della capacità per area geografica rivela una chiara specializzazione strategica e un'allocazione del capitale mirata:
    \vspace{0.2cm}

    \begin{columns}[T]
        \begin{column}{0.48\textwidth}
            \begin{block}{Stati Uniti: Scala Eolica}
                Concentrato oltre il \textbf{50\% della capacità eolica} totale di Gruppo (\textbf{6.218 MW}). Sfruttamento di economie di scala e incentivi federali.
            \end{block}
        \end{column}
        \pause
        \note{Negli Stati Uniti si concentra la scala eolica con oltre seimila megawatt di wind, ottimizzando gli incentivi federali.}
        
        \begin{column}{0.48\textwidth}
            \begin{block}{Italia: Roccaforte Idrica}
                Presidio degli asset storici ad alto margine con \textbf{12.994 MW} di idroelettrico. Base per la flessibilità del mercato europeo.
            \end{block}
        \end{column}
    \end{columns}
    \vspace{0.3cm}
    \pause
    \note{L'Italia si conferma la roccaforte idroelettrica di Gruppo, garantendo la flessibilità necessaria sul mercato europeo.}
    
    \begin{alertblock}{Sintesi per l'Analista}
        La tecnologia produttiva segue la vocazione del territorio: scala e incentivi negli USA, flessibilità regolata e stabilità di margine in Italia.
    \end{alertblock}
    \note{In sintesi, l'allocazione segue la vocazione territoriale: crescita e scala negli Stati Uniti, flessibilità regolata e margini stabili in Italia.}
\end{frame}
